% PVAS — WACV 2027 submission: SUPPLEMENTARY MATERIAL (standalone document).
% Build: pdflatex supplement (no bibtex needed; no \cite in this document —
% all citations live in the main paper).
% Cross-references into the main paper are hardcoded ("Eq. (4) of the main
% paper", "Sec. 3.5", ...) and were verified against main.aux at edit time;
% re-verify if equations/sections are added to or removed from the main paper.
\documentclass[10pt,twocolumn,letterpaper]{article}

\usepackage{wacv}
\wacvalgorithmstrack % WACV 2027 Algorithms track (docs/paper_plan.md)
\usepackage{times}
\usepackage{epsfig}
\usepackage{graphicx}
\usepackage{amsmath,amssymb}
\usepackage{booktabs}
\usepackage{multirow}
\usepackage{algorithm}
\usepackage{algpseudocode}
\usepackage{xcolor}
\usepackage[pagebackref=true,breaklinks=true,colorlinks,bookmarks=false]{hyperref}

% \wacvfinalcopy % KEEP COMMENTED for anonymous submission
\def\wacvPaperID{****} % insert assigned OpenReview ID at submission
\ifwacvfinal\pagestyle{empty}\fi

% Shared macros for the PVAS WACV paper. All sections use ONLY these.
\newcommand{\pvas}{\textsc{pvas}}
\newcommand{\pigdm}{$\Pi$GDM}
\newcommand{\nfe}{\textsc{nfe}}
\newcommand{\psnrm}{\textsc{psnr}}
\newcommand{\ssimm}{\textsc{ssim}}
\newcommand{\xz}{x_0}
\newcommand{\xt}{x_t}
\newcommand{\xtilde}{\tilde{x}_t}
\newcommand{\abar}{\bar{\alpha}_t}
\newcommand{\sigt}{\sigma_t}
\newcommand{\mut}{\mu_\theta}
\newcommand{\PV}{\hat{P}}
\newcommand{\Amat}{A}
\newcommand{\gain}{K}
\newcommand{\RLOCAL}{\textsuperscript{\dag}}   % reproduced locally
\newcommand{\CITED}{\textsuperscript{\S}}      % from cited paper
\newcommand{\PENDING}{---}                     % experiment specified, not run
\newcommand{\TODO}[1]{\textcolor{red}{[#1]}}


% Supplement numbering: sections A, B, ...; floats S1, S2, ...
\renewcommand{\thesection}{\Alph{section}}
\renewcommand{\thetable}{S\arabic{table}}
\renewcommand{\thefigure}{S\arabic{figure}}
\renewcommand{\thealgorithm}{S\arabic{algorithm}}
\renewcommand{\theequation}{S\arabic{equation}}

\begin{document}

\title{Supplementary Material:\\One Trajectory Is Enough}
\author{Anonymous WACV submission\\Paper ID \wacvPaperID}
\maketitle

\noindent This supplement contains material moved out of the main paper
for space: the domain-agnostic \pvas{} loop (Alg.~\ref{alg:pvas}), the
proof of Proposition~1 (Sec.~\ref{sec:supp_proof}),
qualitative CelebA-HQ panels (Fig.~\ref{fig:qual_sr}), a single-slice
MRI PV-map diagnostic (Fig.~\ref{fig:pvmaps_mri}), the tabular form of
the CIFAR-10 budget sweep (Tab.~\ref{tab:sr_budget}), and the pending
MRI grids (Tabs.~\ref{tab:mri80} and~\ref{tab:mri_r8}). Provenance tags
follow the main paper: \RLOCAL{} marks results reproduced with the code
in this supplement on a single Apple M4 Pro (MPS/CPU); \CITED{} marks
numbers taken from the cited paper; \PENDING{} marks experiments that
are fully specified but not yet complete --- no projected or estimated
numbers appear anywhere.

\section{The \pvas{} algorithm}\label{sec:supp_alg}

Algorithm~\ref{alg:pvas} states the domain-agnostic \pvas{} loop
referenced in Secs.~3.2--3.4 of the main paper. Equation numbers below
refer to the main paper: the data-consistency correction is Eq.~(2),
the probe estimator and optional smoothing are Eqs.~(4)--(7), the
acquisition score is Eq.~(9), and the optional per-coordinate gain is
Eq.~(10). Both production instantiations (main paper Sec.~3.4) run
selection-only: $\PV$ is computed on demand at the acquisition rounds
and nowhere else, and the data-consistency gain stays the scalar
\pigdm-style one. D1 (multi-coil MRI): $q{=}2$, $R_a{=}2$, $T{=}20$
($28$ \nfe). D2 (guided super-resolution): $q{=}4$, $R_a{=}3$,
$T{=}100$ ($124$ \nfe). The per-step gate variant (main paper
Sec.~3.3, ablation) additionally refreshes $\PV$ every $k$ steps and
replaces the scalar gain with the per-coordinate one.

\begin{algorithm}[t]
\caption{\pvas: variance-guided active sensing in a single trajectory}
\label{alg:pvas}
\begin{algorithmic}[1]
\Require denoiser $\mut$; schedule $\sigma_T > \dots > \sigma_1$; initial
atoms $M \gets M_0$, measurements $y$, noise $\{\sigma_i^2\}$; budget $L$
over rounds $\mathcal{T}_{\mathrm{acq}} \subset [t_a, t_b]$; probes $q$;
optional gate refresh period $k$ (default: selection-only, no gate)
\State $\xt \gets \sigma_T\,\epsilon$, \quad $\epsilon \sim \mathcal{N}(0, I)$
\For{$t = T, \dots, 1$}
  \State $\hat{\xz} \gets \mut(\xt, t)$ \Comment{1 \nfe}
  \If{$t \in \mathcal{T}_{\mathrm{acq}}$ \textbf{or} (gate variant: $t \equiv 0 \ (\mathrm{mod}\ k)$)}
    \State $\PV \gets \Pi_{\mathcal{F}}\big(\tfrac{\sigt^2}{q} \sum_{m=1}^{q}
      \phi(J_t v^{(m)}, v^{(m)})\big)$ \Comment{$2q$ \nfe; Eqs.~(4)--(7), main paper}
  \EndIf
  \State $\gain_i \gets \sigt^2 / (\sigt^2 + \sigma_i^2)$ \Comment{default; gate variant: $\PV_i / (\PV_i + \sigma_i^2)$, Eq.~(10)}
  \State $\hat{\xz} \gets \hat{\xz} + \tilde{\Amat}_M^{H}\big(\gain \odot (y - \Amat_M \hat{\xz})\big)$ \Comment{Eq.~(2), main paper}
  \If{$t \in \mathcal{T}_{\mathrm{acq}}$}
    \State $M \gets M \cup \{\text{top-}L_r \text{ atoms by } s(a), \text{Eq.~(9), main paper}\}$; extend $y$
    \State new $M$ enters Eq.~(2): same step on the raw $\mut(\xt,t)$ (D2) or at the next step (D1) \Comment{main paper Sec.~3.3}
  \EndIf
  \State $\xt \gets$ DDIM / re-noising step to $\sigma_{t-1}$ from corrected $\hat{\xz}$
\EndFor
\State \Return $\hat{\xz}$
\end{algorithmic}
\end{algorithm}

\section{Proof of Proposition 1}\label{sec:supp_proof}

Proposition~1 (main paper Sec.~3.5) is the law of total variance
instantiated on a mode assignment; we record the four lines for
completeness and claim no depth beyond the instantiation. Let $m$ be a
discrete latent with
$p(\xz \mid y) = \sum_m p(m \mid y)\, p(\xz \mid y, m)$, and work
coordinate-wise. Conditioning the second moment on $m$,
\begin{equation}
  \mathbb{E}[\xz^2 \mid y]
  = \mathbb{E}_{m \mid y}\big[\operatorname{Var}[\xz \mid y, m]
    + \mathbb{E}[\xz \mid y, m]^2\big],
  \label{eq:supp_ltv1}
\end{equation}
and subtracting
$\mathbb{E}[\xz \mid y]^2 = \big(\mathbb{E}_{m \mid y}\,
\mathbb{E}[\xz \mid y, m]\big)^2$ gives
\begin{equation}
  \operatorname{Var}[\xz \mid y]
  = \mathbb{E}_{m \mid y}\big[\operatorname{Var}[\xz \mid y, m]\big]
  + \operatorname{Var}_{m \mid y}\big(\mathbb{E}[\xz \mid y, m]\big).
  \label{eq:supp_ltv2}
\end{equation}
For the second claim, model each of $K$ independent trajectories as
committing to a mode $m^{(j)} \sim p(m \mid y)$, i.i.d., and reporting
the within-mode mean $\hat{\xz}^{(j)} = \mathbb{E}[\xz \mid y, m^{(j)}]$.
The Bessel-corrected sample variance of
$\{\hat{\xz}^{(j)}\}_{j=1}^{K}$ is unbiased for
$\operatorname{Var}_{m \mid y}\big(\mathbb{E}[\xz \mid y, m]\big)$ for
any $K \ge 2$; at $K{=}2$ it reduces to
$\tfrac{1}{2}\big(\hat{\xz}^{(1)} - \hat{\xz}^{(2)}\big)^2$, the
disagreement term of the \pvas-duo score (Eq.~(12), main paper). The
idealization --- trajectories draw modes independently from
$p(m \mid y)$ and report exact within-mode means --- is a model of
trajectory behavior, stated rather than proved; the decomposition
itself is exact.

\section{Qualitative CelebA-HQ results}\label{sec:supp_celeba}

Figure~\ref{fig:qual_sr} shows the full qualitative guided
super-resolution panel on CelebA-HQ $256^2$ referenced in Sec.~4.2 of
the main paper. The protocol mirrors the CIFAR-10 configuration at
CelebA-HQ scale: $8{\times}8$-pixel blocks, $L{=}96$ acquired blocks,
$T{=}100$ steps, $R_a{=}3$ rounds. This panel is illustrative
($n{=}6$ identities); the quantitative CelebA-HQ evaluation is part of
the pending full-scale sweep and no quantitative claim is made from
this figure beyond the per-panel means printed in the caption.

\begin{figure}[t]
  \centering
  \includegraphics[width=\linewidth]{figs/fig_qual_sr.pdf}
  \caption{Qualitative guided SR on CelebA-HQ $256^2$ ($n{=}6$
  identities, $8{\times}8$-pixel blocks, $L{=}96$ acquired blocks,
  $T{=}100$ steps, $R_a{=}3$ rounds)\RLOCAL. Columns: ground truth,
  low-resolution preview, \emph{static}, \emph{random}, \emph{pv}
  (ours), and the PV map with acquired blocks overlaid. Mean \psnrm{}
  on this illustrative panel: static $23.72$, random $23.83$,
  \emph{pv} $25.26$, oracle $25.83$\,dB; the full-scale CelebA-HQ
  quantitative evaluation is pending.}
  \label{fig:qual_sr}
\end{figure}

\section{MRI PV-map diagnostic}\label{sec:supp_pvmaps}

Figure~\ref{fig:pvmaps_mri} visualizes the acquisition signal itself on
a single fastMRI brain slice, under the production D1 configuration of
the main paper (Sec.~3.4): $R{=}8$ base mask, $L{=}16$ active lines,
truncated $\sigma_{\max}{=}10$ schedule with $T{=}20$ EDM steps, and
the production estimator flags --- $q{=}2$ centered Hutchinson probes
in the rank-one positive k-space form, no radial smoothing or
normalization. For this visualization $\PV$ is refreshed at every step;
the production selection-only mode computes it only at the two
acquisition rounds. The three left panels show
$\log_{10}\PV_i(t)$ over k-space at an early, mid, and late step of the
reverse trajectory: the map evolves from a prior-dominated shape early
on toward image-adapted, angularly structured high-variance regions
late in the trajectory --- exactly the structure the per-column score
of Eq.~(10) (main paper) ranks on, and the structure that radial
smoothing would erase. The right panel traces the number of acquired
columns over the trajectory: the mask grows from the base mask in the
two acquisition rounds placed at $50\%$ and $85\%$ of the schedule,
eight columns per round. This figure is a single-slice illustrative
diagnostic of the mechanism; it supports no quantitative claim.

\begin{figure}[t]
  \centering
  \includegraphics[width=\linewidth]{figs/fig_pvmaps_mri.pdf}
  \caption{PV-map diagnostic on one fastMRI brain slice ($R{=}8$,
  $L{=}16$, $\sigma_{\max}{=}10$, $T{=}20$ EDM steps; production D1
  estimator flags: $q{=}2$ centered probes, rank-one positive k-space
  form, no smoothing; per-step refresh for visualization only)\RLOCAL.
  Left three panels:
  $\log_{10}\PV_i(t)$ at an early, mid, and late step. Right:
  acquired-column count over the trajectory (two rounds, at $50\%$ and
  $85\%$ of the schedule). Illustrative single-slice diagnostic; no
  quantitative claim.}
  \label{fig:pvmaps_mri}
\end{figure}

\section{Budget sweep: tabular form}\label{sec:supp_budget}

Table~\ref{tab:sr_budget} gives the numeric values behind the CIFAR-10
budget-sweep curve (Fig.~2 of the main paper): \psnrm{} at matched
budget for $L \in \{6, 12, 18, 24\}$ acquired blocks. The paired
margins and $p$-values are quoted in main paper Sec.~4.2.

% supplement candidate -- not \input by the main paper; kept generated for the supplement (with fig_nfe_cost.pdf).
% Auto-generated by paper/pvas_v2/tables/gen_tables.py -- do not edit.
% input: outputs/poc_natural/cifar_sr_b6/results.json [R-LOCAL]
% input: outputs/poc_natural/cifar_sr_main/results.json [R-LOCAL]
% input: outputs/poc_natural/cifar_sr_b18/results.json [R-LOCAL]
% input: outputs/poc_natural/cifar_sr_b24/results.json [R-LOCAL]
% n = 64/128 images per budget point (from JSON configs)
\begin{table}[t]
  \centering\small
  \setlength{\tabcolsep}{4pt}
  \caption{\psnrm{} (dB) vs.\ acquisition budget $L$ (blocks) on CIFAR-10 adaptive super-resolution, $n{=}64/128$ images per point\RLOCAL. Best non-oracle entry per row in bold; the oracle column is a reference. See also Fig.~\ref{fig:budget}.}
  \label{tab:sr_budget}
  \begin{tabular}{lcccc}
    \toprule
    Budget $L$ & Static & Random & \pvas{} (ours) & Oracle (ref.) \\
    \midrule
    6 & 19.34\,$\pm$\,2.35 & 19.50\,$\pm$\,2.35 & \textbf{20.18\,$\pm$\,2.72} & 21.12\,$\pm$\,2.73 \\
    12 & 19.79\,$\pm$\,2.50 & 19.98\,$\pm$\,2.46 & \textbf{21.43\,$\pm$\,3.70} & 22.95\,$\pm$\,3.80 \\
    18 & 20.70\,$\pm$\,2.42 & 20.80\,$\pm$\,2.43 & \textbf{23.29\,$\pm$\,4.00} & 25.40\,$\pm$\,4.50 \\
    24 & 21.69\,$\pm$\,2.45 & 21.76\,$\pm$\,2.40 & \textbf{25.56\,$\pm$\,6.22} & 27.85\,$\pm$\,6.06 \\
    \bottomrule
  \end{tabular}
\end{table}
     % -> tab:sr_budget

\section{Additional MRI configurations}\label{sec:supp_pending}

The $R{=}8$, $L{=}16$ configuration completed and moved to Table~6 of
the main paper. The $R{=}8$, $L{=}8$ configuration
(Tab.~\ref{tab:mri_r8}, $n{=}30$)\RLOCAL{} shows the same ordering at
half the budget. The $\sigma_{\max}{=}80$ (from-scratch) grid
(Tab.~\ref{tab:mri80}, $n{=}30$ per configuration)\RLOCAL{} uses the
same protocol, methods, and provenance rules as Table~6 of the main
paper. Margins are larger than at $\sigma_{\max}{=}10$: \pvas{} gains
$+0.92\pm0.44$\,dB over the matched static mask at $R{=}4$, $L{=}8$ and
$+2.15\pm0.50$\,dB at $R{=}8$, $L{=}16$ (both
$p{=}9.3{\times}10^{-10}$), and in this regime exceeds the greedy
residual-oracle policy ($+0.11$\,dB, $p{=}0.014$ and $+0.53$\,dB,
$p{=}1.8{\times}10^{-8}$ respectively) --- greedy error-chasing is not
the optimal policy, and the variance score selects better when the
trajectory starts from an uninformed initialization. Two caveats: the
passive $\sigma_{\max}{=}10$ baseline is stronger in absolute terms at
$R{=}4$ (Table~6, main paper), so part of this margin reflects
rescuing a weaker-initialized sampler; and the gate-only row remains
below \pigdm{} here as well.

% Auto-generated by paper/pvas_v2/tables/gen_tables.py -- do not edit.
% input: outputs/PVAS_ng80/R4_a8/results.json [MISSING -> \PENDING]
% input: outputs/PVAS_ng80/R8_a16/results.json [MISSING -> \PENDING]
% slices per config: R4_a8: n=PENDING; R8_a16: n=PENDING
\begin{table}[t]
  \centering\small
  \setlength{\tabcolsep}{3pt}
  \caption{Multi-coil fastMRI brain k-space line acquisition ($\sigma_{\max}{=}80$ regime, selection-only). All adaptive methods receive the same total column budget as the static baseline (matched budget); SENSE-referenced \psnrm/\ssimm. Best non-oracle entry per group in bold. $\Delta$ rows: paired per-slice \psnrm{} differences with one-sided Wilcoxon $p$.}
  \label{tab:mri80}
  \begin{tabular}{lcccc}
    \toprule
    Method & Cols & \psnrm{} (dB) & \ssimm{} & \nfe/slice \\
    \midrule
    \multicolumn{5}{l}{\emph{$R{=}4$, $L{=}8$ active lines}} \\
    \pigdm{} (passive) & \PENDING & \PENDING & \PENDING & \PENDING \\
    Static, matched budget & \PENDING & \PENDING & \PENDING & \PENDING \\
    Random adaptive & \PENDING & \PENDING & \PENDING & \PENDING \\
    \pvas{} (ours) & \PENDING & \PENDING & \PENDING & \PENDING \\
    Ensemble-$N$ particles (ADS-style ref.) & \PENDING & \PENDING & \PENDING & \PENDING \\
    Oracle selection (bound) & \PENDING & \PENDING & \PENDING & \PENDING \\
    \midrule
    \multicolumn{5}{l}{\emph{$R{=}8$, $L{=}16$ active lines}} \\
    \pigdm{} (passive) & \PENDING & \PENDING & \PENDING & \PENDING \\
    Static, matched budget & \PENDING & \PENDING & \PENDING & \PENDING \\
    Random adaptive & \PENDING & \PENDING & \PENDING & \PENDING \\
    \pvas{} (ours) & \PENDING & \PENDING & \PENDING & \PENDING \\
    Ensemble-$N$ particles (ADS-style ref.) & \PENDING & \PENDING & \PENDING & \PENDING \\
    Oracle selection (bound) & \PENDING & \PENDING & \PENDING & \PENDING \\
    \bottomrule
  \end{tabular}
\end{table}
         % -> tab:mri80

\begin{table}[t]
  \centering\small
  \setlength{\tabcolsep}{3pt}
  \caption{Multi-coil fastMRI brain k-space line acquisition, $R{=}8$,
  $L{=}8$ configuration of the $\sigma_{\max}{=}10$ regime ($T{=}20$
  EDM steps, $n{=}30$)\RLOCAL. Protocol identical to Table~6 of the
  main paper: matched total column budget, SENSE-referenced
  \psnrm/\ssimm; $\Delta$ rows are paired per-slice differences with
  one-sided Wilcoxon $p$.}
  \label{tab:mri_r8}
  \begin{tabular}{lcccc}
    \toprule
    Method & Cols & \psnrm{} (dB) & \ssimm{} & \nfe/slice \\
    \midrule
    \multicolumn{5}{l}{\emph{$R{=}8$, $L{=}8$ active lines}} \\
    Static, matched budget & 51 & 27.92\,$\pm$\,1.15 & 0.781\,$\pm$\,0.038 & 20 \\
    Random adaptive & 51 & 27.96\,$\pm$\,1.17 & 0.785\,$\pm$\,0.039 & 100 \\
    \pvas{} (ours) & 51 & \textbf{28.47\,$\pm$\,1.18} & \textbf{0.795\,$\pm$\,0.039} & 100 \\
    $\Delta$(\pvas{} $-$ static match) & ($n{=}30$) & $+0.55$\,$\pm$\,$0.28$ & $p{=}9.3{\times}10^{-10}$ & \\
    $\Delta$(\pvas{} $-$ random adaptive) & ($n{=}30$) & $+0.52$\,$\pm$\,$0.23$ & $p{=}9.3{\times}10^{-10}$ & \\
    \bottomrule
  \end{tabular}
\end{table}

\end{document}
